
\title{\textbf{Mathematical Foundations, Architecture Diagrams, and State Analysis\\
of Modern LLM Training Methods --- Part 3:\\
Data Generation, Curation, Distillation, Compression,\\
Agentic, Multimodal, Knowledge Injection, and Deployment}}
\author{Chief Strategist J}
\date{\today}
\maketitle
\tableofcontents
\newpage

%% ============================================================
\part{Synthetic Data Generation}
%% ============================================================

\section{Self-Instruct}

\subsection{Mathematical Foundation}
Self-Instruct bootstraps an instruction dataset from a small seed pool $\mathcal{S}$ by using the LLM itself as a generator. At each iteration, $k$ seed instructions are sampled and the model generates a new instruction $\hat{x}$, its input, and output:
\begin{equation}
  (\hat{x}, \hat{i}, \hat{y}) \sim P_\theta(\cdot \mid \text{few-shot}(\mathcal{S}))
\end{equation}
Generated examples pass a quality filter (ROUGE similarity $< \delta$ to existing pool, length check):
\begin{equation}
  \mathcal{S} \leftarrow \mathcal{S} \cup \{(\hat{x},\hat{i},\hat{y})\} \;\text{ iff }\; \max_{s\in\mathcal{S}} \text{ROUGE}(\hat{x}, s) < \delta
\end{equation}
The growing pool is used to fine-tune $\theta$, improving the quality of subsequent generations.

\subsection{Architecture Flow Diagram}

\begin{figure}[h]
\centering
\begin{tikzpicture}[node distance=0.7cm and 1.0cm]
  \node[datblk] (seed) {Seed Pool\\$\mathcal{S}$ (175 seeds)};
  \node[grayblk, right=of seed] (sample) {Sample $k$\\Instructions};
  \node[trnblk, right=of sample] (gen) {LLM\\Generates\\$(\hat{x},\hat{i},\hat{y})$};
  \node[orangeblk, right=of gen] (filt) {Quality\\Filter\\ROUGE $< \delta$};
  \node[datblk, right=of filt] (pool) {Expanded\\Pool $\mathcal{S}'$};
  \node[mrgblk, below=0.9cm of pool] (sft) {SFT on\\$\mathcal{S}'$};
  \node[trnblk, left=2.6cm of sft] (better) {Better\\LLM $\theta'$};

  \draw[arr] (seed)--(sample);
  \draw[arr] (sample)--(gen);
  \draw[arr] (gen)--(filt);
  \draw[arr] (filt)--(pool);
  \draw[arr] (pool)--(sft);
  \draw[arr] (sft)--(better);
  \draw[arr,dashed,green!60!black] (better.north) |- (gen.south) node[midway,below,font=\tiny]{iterate};
  \draw[arr,dashed,blue!60] (pool.north) -- ++(0,0.4) -| (seed.north) node[midway,above,font=\tiny]{pool grows};
\end{tikzpicture}
\caption{Self-Instruct: LLM generates new instructions from seed pool; quality filter expands the pool for SFT.}
\end{figure}

\subsection{State Transition Table}
\begin{table}[h]\centering
\begin{tabular}{lll}\toprule
\textbf{Property} & \textbf{Before} & \textbf{After}\\\midrule
Pool size & 175 seed instructions & $52{,}000+$ generated\\
Annotation source & Human-written seeds & LLM-generated\\
Diversity filter & N/A & ROUGE deduplication\\
Training data cost & Low (seeds only) & Essentially free to scale\\
Instruction following & Weak & Strong\\\bottomrule
\end{tabular}
\caption{State properties before and after Self-Instruct data generation.}
\end{table}

\newpage
%% ------------------------------------------------------------
\section{Evol-Instruct}

\subsection{Mathematical Foundation}
Evol-Instruct evolves existing instructions into harder variants via two mutation strategies. \emph{In-depth} evolution adds constraints, increases steps, or requires deeper reasoning. \emph{In-breadth} evolution creates new topic-adjacent instructions:
\begin{equation}
  \hat{x} = M_\text{depth}(x) \;\text{ or }\; \hat{x} = M_\text{breadth}(x), \quad \hat{x} \sim P_{\text{teacher}}(\cdot \mid \text{evolve-prompt}, x)
\end{equation}
Evolved instructions are validated (non-empty, no refusal, dissimilar from parent) and collected into $\mathcal{D}_{\text{evol}}$. Difficulty increases monotonically with evolution rounds $E$:
\begin{equation}
  \mathbb{E}[\text{difficulty}(\hat{x}_E)] > \mathbb{E}[\text{difficulty}(\hat{x}_{E-1})]
\end{equation}

\subsection{Architecture Flow Diagram}

\begin{figure}[h]
\centering
\begin{tikzpicture}[node distance=0.6cm and 1.0cm]
  \node[datblk] (orig) {Original\\Instruction $x$};
  \node[trnblk, above right=0.5cm and 1.2cm of orig] (depth) {In-Depth\\Mutation};
  \node[trnblk, below right=0.5cm and 1.2cm of orig] (breadth) {In-Breadth\\Mutation};
  \node[grayblk, right=3.0cm of orig] (evolved) {Evolved\\$\hat{x}$};
  \node[orangeblk, right=of evolved] (val) {Validate\\(reject refusals)};
  \node[datblk, right=of val] (pool) {Evolved\\Dataset};
  \node[mrgblk, below=0.8cm of pool] (sft) {SFT\\(WizardLM style)};

  \draw[arr] (orig.east) |- (depth.west);
  \draw[arr] (orig.east) |- (breadth.west);
  \draw[arr] (depth.east) -| (evolved.north);
  \draw[arr] (breadth.east) -| (evolved.south);
  \draw[arr] (evolved)--(val);
  \draw[arr] (val)--(pool);
  \draw[arr] (pool)--(sft);
  \draw[arr,dashed,green!60!black] (pool.north) -- ++(0,0.5) -| (orig.north) node[midway,above,font=\tiny]{$E$ rounds};
\end{tikzpicture}
\caption{Evol-Instruct: in-depth and in-breadth mutations generate harder variants; validated pool used for SFT.}
\end{figure}

\subsection{State Transition Table}
\begin{table}[h]\centering
\begin{tabular}{lll}\toprule
\textbf{Property} & \textbf{Before} & \textbf{After}\\\midrule
Instruction difficulty & Low (seed) & Progressively higher (per round $E$)\\
Instruction diversity & Limited & Broad via breadth mutations\\
Data source & Human-written & LLM-evolved\\
Evolution rounds & 0 & $E$ (typically 4--8)\\
Model performance & Baseline & WizardLM-style improvements\\\bottomrule
\end{tabular}
\caption{State properties before and after Evol-Instruct.}
\end{table}

\newpage
%% ------------------------------------------------------------
\section{Evol-Code}

\subsection{Mathematical Foundation}
Evol-Code applies the same evolutionary principle to programming tasks, using code-specific mutation operators: adding input/output constraints, combining two functions, increasing time/space complexity, or switching programming language:
\begin{equation}
  \hat{p} = M_{\text{code}}(p),\quad M \in \{\text{constrain}, \text{combine}, \text{optimize}, \text{translate}\}
\end{equation}
Each evolved problem $\hat{p}$ is solved by a teacher model, and the (problem, solution) pair is validated by unit-test execution:
\begin{equation}
  (\hat{p},\hat{s}) \in \mathcal{D}_{\text{code}} \iff \text{unit\_tests}(\hat{s}, \hat{p}) = \text{PASS}
\end{equation}

\subsection{Architecture Flow Diagram}

\begin{figure}[h]
\centering
\begin{tikzpicture}[node distance=0.6cm and 1.0cm]
  \node[datblk] (seed) {Seed Code\\Problems};
  \node[trnblk, right=of seed] (mutate) {Code\\Mutator};
  \node[grayblk, right=of mutate] (prob) {Evolved\\Problem $\hat{p}$};
  \node[orangeblk, right=of prob] (solve) {Teacher\\LLM Solves};
  \node[grayblk, right=of solve] (sol) {Solution $\hat{s}$};
  \node[cyanblk, below=0.9cm of sol] (test) {Unit Test\\Execution};
  \node[datblk, left=1.4cm of test] (pass) {Pass: add to\\$\mathcal{D}_{\text{code}}$};
  \node[lossblk, right=0.8cm of test] (fail) {Fail: discard};

  \draw[arr] (seed)--(mutate);
  \draw[arr] (mutate)--(prob);
  \draw[arr] (prob)--(solve);
  \draw[arr] (solve)--(sol);
  \draw[arr] (sol)--(test);
  \draw[arr] (test)--(pass);
  \draw[arr] (test)--(fail);
  \draw[arr,dashed,green!60!black] (pass.west) -- ++(-0.5,0) |- (seed.west) node[midway,left,font=\tiny]{grow pool};
\end{tikzpicture}
\caption{Evol-Code: code problems are mutated, solved by teacher, validated by unit tests before dataset inclusion.}
\end{figure}

\subsection{State Transition Table}
\begin{table}[h]\centering
\begin{tabular}{lll}\toprule
\textbf{Property} & \textbf{Before} & \textbf{After}\\\midrule
Problem complexity & Simple seed tasks & Multi-constraint, optimised\\
Validation method & Human review & Automated unit tests\\
Data quality gate & None & Unit-test pass rate\\
Programming languages & Fixed & Multi-language via translate mutation\\
Code reasoning ability & Baseline & Significantly improved\\\bottomrule
\end{tabular}
\caption{State properties before and after Evol-Code data generation.}
\end{table}

\newpage
%% ------------------------------------------------------------
\section{Synthetic Data Generation}

\subsection{Mathematical Foundation}
A powerful teacher model $P_T$ generates synthetic (input, output) pairs from topic/schema prompts, which are then used to train a smaller student $P_S$. The synthetic data distribution approximates the real task distribution $P_{\text{real}}$:
\begin{equation}
  \mathcal{D}_{\text{synth}} = \{(x_i, y_i) : x_i \sim P_{\text{topic}},\; y_i \sim P_T(\cdot \mid x_i)\}
\end{equation}
Quality is improved by filtering on a scoring model or perplexity threshold $\tau_p$:
\begin{equation}
  (x,y) \in \mathcal{D}_{\text{final}} \iff \text{score}(x,y) \geq \tau_s \;\text{ and }\; \text{PPL}(y) \leq \tau_p
\end{equation}

\subsection{Architecture Flow Diagram}

\begin{figure}[h]
\centering
\begin{tikzpicture}[node distance=0.7cm and 1.0cm]
  \node[datblk] (schema) {Schema /\\Topic Seeds};
  \node[frzblk, right=of schema] (teacher) {Teacher Model\\$P_T$ (large)};
  \node[grayblk, right=of teacher] (synth) {Synthetic\\Pairs $(x,y)$};
  \node[orangeblk, right=of synth] (score) {Quality\\Scorer};
  \node[datblk, right=of score] (final) {Filtered\\$\mathcal{D}_{\text{final}}$};
  \node[trnblk, below=0.9cm of final] (student) {Train Student\\$P_S$ (small)};

  \draw[arr] (schema)--(teacher);
  \draw[arr] (teacher)--(synth);
  \draw[arr] (synth)--(score);
  \draw[arr] (score)--(final);
  \draw[arr] (final)--(student);
  \node[red,font=\tiny,below=0.2cm of score] {reject low-quality};
\end{tikzpicture}
\caption{Synthetic Data Generation: teacher generates pairs from schema seeds; quality filter builds training set for student.}
\end{figure}

\subsection{State Transition Table}
\begin{table}[h]\centering
\begin{tabular}{lll}\toprule
\textbf{Property} & \textbf{Before} & \textbf{After}\\\midrule
Data source & Human annotation & LLM-generated at scale\\
Annotation cost & High & Very low (API cost)\\
Data volume & Thousands & Millions scalable\\
Quality control & Human review & Automated scoring\\
Domain coverage & Limited & Any topic via schema\\\bottomrule
\end{tabular}
\caption{State properties before and after synthetic data generation.}
\end{table}

\newpage
%% ------------------------------------------------------------
\section{Rejection Sampling}

\subsection{Mathematical Foundation}
Rejection Sampling generates $N$ candidate completions from the current policy and keeps only those that pass a quality threshold, building a filtered fine-tuning corpus:
\begin{equation}
  \mathcal{D}_{\text{RS}} = \{(x, y^{(k)}) : y^{(k)} \sim \pi_\theta(x),\; r(x, y^{(k)}) \geq \tau,\; k = 1..N\}
\end{equation}
Expected acceptance rate at threshold $\tau$: $\alpha = P_{y \sim \pi_\theta}[r(x,y) \geq \tau]$. After SFT on $\mathcal{D}_{\text{RS}}$, the model's probability mass shifts toward high-reward completions:
\begin{equation}
  \mathbb{E}_{y \sim \pi_{\theta'}}[r(x,y)] > \mathbb{E}_{y \sim \pi_\theta}[r(x,y)]
\end{equation}

\subsection{Architecture Flow Diagram}

\begin{figure}[h]
\centering
\begin{tikzpicture}[node distance=0.7cm and 1.0cm]
  \node[datblk] (x) {Prompts $\{x_i\}$};
  \node[trnblk, right=of x] (pol) {Policy $\pi_\theta$\\Sample $N$};
  \node[grayblk, right=of pol, minimum height=2.0cm] (cands) {$N$ Candidates\\per prompt};
  \node[orangeblk, right=of cands] (rm) {Reward Model\\$r(x,y)$};
  \node[datblk, right=of rm] (keep) {Keep\\$r \geq \tau$};
  \node[lossblk, below=0.9cm of cands] (drop) {Discard\\$r < \tau$};
  \node[mrgblk, right=0.8cm of keep] (sft) {SFT on\\accepted};

  \draw[arr] (x)--(pol);
  \draw[arr] (pol)--(cands);
  \draw[arr] (cands)--(rm);
  \draw[arr] (rm)--(keep);
  \draw[arr] (rm.south) -- ++(0,-0.3) -| (drop.east);
  \draw[arr] (keep)--(sft);
\end{tikzpicture}
\caption{Rejection Sampling: $N$ completions generated; only those above reward threshold kept for SFT.}
\end{figure}

\subsection{State Transition Table}
\begin{table}[h]\centering
\begin{tabular}{lll}\toprule
\textbf{Property} & \textbf{Before} & \textbf{After}\\\midrule
Training data quality & Mixed & High (above threshold $\tau$)\\
Samples per prompt & 1 & $N$ (with $\alpha N$ accepted)\\
Reward model & Optional & Required as filter\\
Iteration benefit & None & Compound with model updates\\
RL complexity & None & None (SFT only)\\\bottomrule
\end{tabular}
\caption{State properties before and after Rejection Sampling.}
\end{table}

\newpage
%% ------------------------------------------------------------
\section{Best-of-N Sampling}

\subsection{Mathematical Foundation}
Best-of-N (BoN) sampling selects the highest-reward completion from $N$ independent samples at inference time:
\begin{equation}
  y^* = \arg\max_{k \in 1..N} r(x, y^{(k)}),\quad y^{(k)} \sim \pi_\theta(\cdot \mid x)
\end{equation}
The expected reward improvement over single-sample baseline scales as:
\begin{equation}
  \mathbb{E}[\max_k r(x,y^{(k)})] \approx \mu_r + \sigma_r \cdot \Phi^{-1}\!\left(\frac{N}{N+1}\right)
\end{equation}
where $\mu_r, \sigma_r$ are the reward mean and std under $\pi_\theta$, and $\Phi^{-1}$ is the inverse normal CDF. BoN can also be used to build SFT datasets by collecting $(x, y^*)$ pairs.

\subsection{Architecture Flow Diagram}

\begin{figure}[h]
\centering
\begin{tikzpicture}[node distance=0.5cm and 0.9cm]
  \node[datblk,minimum width=1.8cm] (x) {Prompt $x$};
  \node[trnblk,minimum width=1.6cm, above right=0.3cm and 1.0cm of x] (y1) {$y^{(1)}$};
  \node[trnblk,minimum width=1.6cm, right=1.0cm of x] (y2) {$y^{(2)}$};
  \node[trnblk,minimum width=1.6cm, below right=0.3cm and 1.0cm of x] (yN) {$y^{(N)}$};
  \node[orangeblk, right=1.6cm of y2] (rm) {Reward\\$r(x,y^{(k)})$};
  \node[mrgblk, right=of rm] (best) {$y^*=\arg\max$};

  \draw[arr] (x.east) |- (y1.west);
  \draw[arr] (x)--(y2);
  \draw[arr] (x.east) |- (yN.west);
  \draw[arr] (y1.east) -| (rm.north);
  \draw[arr] (y2)--(rm);
  \draw[arr] (yN.east) -| (rm.south);
  \draw[arr] (rm)--(best);

  \node[font=\tiny,below=0.2cm of y2] {$N$ samples};
\end{tikzpicture}
\caption{Best-of-N: $N$ samples scored by reward model; highest reward selected as final output.}
\end{figure}

\subsection{State Transition Table}
\begin{table}[h]\centering
\begin{tabular}{lll}\toprule
\textbf{Property} & \textbf{Before (single)} & \textbf{After (BoN)}\\\midrule
Samples per query & 1 & $N$\\
Compute at inference & $1\times$ & $N\times$\\
Expected reward & $\mu_r$ & $\mu_r + O(\sigma_r \log N)$\\
Model weights & Unchanged & Unchanged\\
Best usage & Cheap inference & Quality-critical generation\\\bottomrule
\end{tabular}
\caption{State properties comparing single-sample and Best-of-N sampling.}
\end{table}

\newpage
%% ============================================================
\part{Data Curation}
%% ============================================================

\section{Data Filtering}

\subsection{Mathematical Foundation}
Data filtering removes low-quality examples from a raw corpus using a set of quality signals $\{q_1,\ldots,q_M\}$. A sample $(x,y)$ is accepted iff it passes all signal thresholds:
\begin{equation}
  (x,y) \in \mathcal{D}_{\text{clean}} \iff \bigwedge_{m=1}^M q_m(x,y) \geq \tau_m
\end{equation}
Common signals include perplexity under a small LM (detects gibberish), toxic content classifiers, length ratios, and language identification. A learned quality classifier $c_\phi$ can consolidate multiple signals:
\begin{equation}
  c_\phi(x,y) = \sigma(w^\top [q_1(x,y),\ldots,q_M(x,y)])
\end{equation}

\subsection{Architecture Flow Diagram}

\begin{figure}[h]
\centering
\begin{tikzpicture}[node distance=0.6cm and 1.0cm]
  \node[datblk] (raw) {Raw Corpus\\(noisy, large)};
  \node[grayblk, right=of raw] (ppl) {Perplexity\\Filter};
  \node[grayblk, right=of ppl] (tox) {Toxicity\\Classifier};
  \node[grayblk, right=of tox] (lang) {Language\\ID};
  \node[orangeblk, right=of lang] (thresh) {Threshold\\Gates $\tau_m$};
  \node[trnblk, above right=0.4cm and 1.0cm of thresh] (pass) {PASS\\$\rightarrow$ keep};
  \node[lossblk, below right=0.4cm and 1.0cm of thresh] (fail) {FAIL\\$\rightarrow$ drop};
  \node[datblk, right=1.8cm of thresh] (clean) {Filtered\\$\mathcal{D}_{\text{clean}}$};

  \draw[arr] (raw)--(ppl);
  \draw[arr] (ppl)--(tox);
  \draw[arr] (tox)--(lang);
  \draw[arr] (lang)--(thresh);
  \draw[arr] (thresh.east) |- (pass.west);
  \draw[arr] (thresh.east) |- (fail.west);
  \draw[arr] (pass.east) -| (clean.north);
\end{tikzpicture}
\caption{Data filtering pipeline: multiple quality signals applied in sequence; only passing samples retained.}
\end{figure}

\subsection{State Transition Table}
\begin{table}[h]\centering
\begin{tabular}{lll}\toprule
\textbf{Property} & \textbf{Before} & \textbf{After}\\\midrule
Corpus size & Very large (noisy) & Smaller (high quality)\\
Noise level & High & Low\\
Domain coverage & Broad (noisy) & Focused (clean)\\
Filter signals & None & $M$ quality signals\\
Training efficiency & Low & Higher (less wasted compute)\\\bottomrule
\end{tabular}
\caption{State properties before and after data filtering.}
\end{table}

\newpage
%% ------------------------------------------------------------
\section{Deduplication}

\subsection{Mathematical Foundation}
Deduplication removes near-duplicate documents to prevent memorisation and improve sample diversity. MinHash LSH estimates Jaccard similarity between documents:
\begin{equation}
  J(A,B) = \frac{|A \cap B|}{|A \cup B|} \approx \frac{1}{K}\sum_{k=1}^K \mathbf{1}[h_k(A) = h_k(B)]
\end{equation}
where $\{h_k\}$ are $K$ independent hash functions and $h_k(A) = \min_{w \in \text{ngrams}(A)} h_k(w)$. Documents with $J(A,B) > \tau_J$ are considered duplicates and one is removed:
\begin{equation}
  \mathcal{D}_{\text{dedup}} = \{d_i : \forall j < i,\; J(d_i, d_j) \leq \tau_J\}
\end{equation}

\subsection{Architecture Flow Diagram}

\begin{figure}[h]
\centering
\begin{tikzpicture}[node distance=0.7cm and 1.0cm]
  \node[datblk] (corpus) {Corpus\\$\mathcal{D}$};
  \node[grayblk, right=of corpus] (ngrams) {N-gram\\Shingles};
  \node[trnblk, right=of ngrams] (minhash) {MinHash\\Signatures};
  \node[grayblk, right=of minhash] (lsh) {LSH\\Bucketing};
  \node[orangeblk, right=of lsh] (cmp) {Pairwise\\$J > \tau_J$?};
  \node[datblk, above right=0.4cm and 1.0cm of cmp] (keep) {Keep (unique)};
  \node[lossblk, below right=0.4cm and 1.0cm of cmp] (drop) {Drop (dup.)};

  \draw[arr] (corpus)--(ngrams);
  \draw[arr] (ngrams)--(minhash);
  \draw[arr] (minhash)--(lsh);
  \draw[arr] (lsh)--(cmp);
  \draw[arr] (cmp.east) |- (keep.west);
  \draw[arr] (cmp.east) |- (drop.west);
\end{tikzpicture}
\caption{Deduplication: MinHash signatures group near-duplicates via LSH; high-similarity pairs discarded.}
\end{figure}

\subsection{State Transition Table}
\begin{table}[h]\centering
\begin{tabular}{lll}\toprule
\textbf{Property} & \textbf{Before} & \textbf{After}\\\midrule
Near-duplicate rate & 15--30\% (web data) & Near zero\\
Memorisation risk & High & Reduced\\
Effective token diversity & Low & High\\
Corpus size & $|\mathcal{D}|$ & $|\mathcal{D}_{\text{dedup}}| < |\mathcal{D}|$\\
Compute wasted on duplicates & High & Eliminated\\\bottomrule
\end{tabular}
\caption{State properties before and after deduplication.}
\end{table}

\newpage
%% ------------------------------------------------------------
\section{Active Learning}

\subsection{Mathematical Foundation}
Active Learning selects the most \emph{informative} unlabelled examples for annotation, maximising model improvement per annotation dollar. Uncertainty sampling selects $x^*$ with highest predictive entropy:
\begin{equation}
  x^* = \arg\max_{x \in \mathcal{U}} \mathbb{H}[P_\theta(y \mid x)] = \arg\max_x -\sum_y P_\theta(y|x)\log P_\theta(y|x)
\end{equation}
Query-by-committee uses disagreement between an ensemble $\{\theta_1,\ldots,\theta_C\}$:
\begin{equation}
  x^* = \arg\max_{x \in \mathcal{U}} \frac{1}{C}\sum_c \mathbb{KL}[P_{\theta_c}(\cdot|x) \| \bar{P}(\cdot|x)]
\end{equation}

\subsection{Architecture Flow Diagram}

\begin{figure}[h]
\centering
\begin{tikzpicture}[node distance=0.7cm and 1.0cm]
  \node[datblk] (unlabeled) {Unlabelled\\Pool $\mathcal{U}$};
  \node[frzblk, right=of unlabeled] (model) {Current\\Model $\theta$};
  \node[orangeblk, right=of model] (unc) {Uncertainty\\Estimator};
  \node[grayblk, right=of unc] (select) {Select top-$b$\\examples};
  \node[datblk, right=of select] (annot) {Human\\Annotator};
  \node[trnblk, below=0.9cm of annot] (retrain) {Retrain\\on $\mathcal{L}$};
  \node[mrgblk, left=2.4cm of retrain] (better) {Better\\$\theta'$};

  \draw[arr] (unlabeled)--(model);
  \draw[arr] (model)--(unc);
  \draw[arr] (unc)--(select);
  \draw[arr] (select)--(annot);
  \draw[arr] (annot)--(retrain);
  \draw[arr] (retrain)--(better);
  \draw[arr,dashed,green!60!black] (better.north) |- (model.south) node[midway,below,font=\tiny]{loop};
\end{tikzpicture}
\caption{Active Learning: uncertainty estimator selects most informative examples; human annotates; model retrains.}
\end{figure}

\subsection{State Transition Table}
\begin{table}[h]\centering
\begin{tabular}{lll}\toprule
\textbf{Property} & \textbf{Before} & \textbf{After}\\\midrule
Annotation strategy & Random sampling & Uncertainty-driven\\
Labels per accuracy gain & Many & Few (targeted)\\
Unlabelled pool & Fully unused & Iteratively queried\\
Annotation budget & Fixed & Allocated by informativeness\\
Model improvement per label & Low & High\\\bottomrule
\end{tabular}
\caption{State properties before and after Active Learning.}
\end{table}

\newpage
%% ------------------------------------------------------------
\section{Curriculum Learning}

\subsection{Mathematical Foundation}
Curriculum Learning trains on examples ordered by difficulty $d(x,y)$, starting with easy examples and progressively introducing harder ones. A competence function $c(t) \in [0,1]$ controls what fraction of the difficulty range is accessible at step $t$:
\begin{equation}
  c(t) = \min\!\left(1,\; c_0 + (1-c_0)\frac{t}{T_{\text{ramp}}}\right)
\end{equation}
The training distribution at step $t$ only samples from examples with $d(x,y) \leq c(t) \cdot d_{\max}$:
\begin{equation}
  \mathcal{D}_t = \{(x,y) : d(x,y) \leq c(t) \cdot d_{\max}\},\quad (x,y) \sim \text{Uniform}(\mathcal{D}_t)
\end{equation}

\subsection{Architecture Flow Diagram}

\begin{figure}[h]
\centering
\begin{tikzpicture}[node distance=0.6cm and 1.0cm]
  \node[datblk] (all) {Full Dataset\\$\mathcal{D}$};
  \node[grayblk, right=of all] (score) {Difficulty\\Scorer $d(\cdot)$};
  \node[orangeblk, right=of score] (sort) {Sort by\\Difficulty};
  \node[trnblk, right=of sort] (easy) {Easy\\Batch $t{=}0$};
  \node[frzblk, below=0.5cm of easy] (med) {Medium\\Batch $t{=}T/2$};
  \node[mrgblk, below=0.5cm of med] (hard) {Hard\\Batch $t{=}T$};
  \node[lossblk, right=1.6cm of med] (model) {Model\\Training};

  \draw[arr] (all)--(score);
  \draw[arr] (score)--(sort);
  \draw[arr] (sort)--(easy);
  \draw[arr] (easy.south)--(med.north);
  \draw[arr] (med.south)--(hard.north);
  \draw[arr] (easy.east) -| (model.north);
  \draw[arr] (med)--(model);
  \draw[arr] (hard.east) -| (model.south);

  \node[font=\tiny,left=0.1cm of easy] {$c(t){=}0.2$};
  \node[font=\tiny,left=0.1cm of med] {$c(t){=}0.6$};
  \node[font=\tiny,left=0.1cm of hard] {$c(t){=}1.0$};
\end{tikzpicture}
\caption{Curriculum Learning: examples sorted by difficulty; competence function progressively unlocks harder batches.}
\end{figure}

\subsection{State Transition Table}
\begin{table}[h]\centering
\begin{tabular}{lll}\toprule
\textbf{Property} & \textbf{Before} & \textbf{After}\\\midrule
Training order & Random & Difficulty-ordered\\
Early training batches & Mixed difficulty & Easy only\\
Convergence speed & Slower & Faster (easier start)\\
Final performance & Baseline & Improved\\
Difficulty metric & N/A & Perplexity, length, or task-specific\\\bottomrule
\end{tabular}
\caption{State properties before and after Curriculum Learning.}
\end{table}

\newpage
%% ============================================================
\part{Knowledge Distillation}
%% ============================================================

\section{Knowledge Distillation}

\subsection{Mathematical Foundation}
Knowledge Distillation transfers the ``dark knowledge'' of a large teacher $P_T$ to a small student $P_S$ by minimising a combination of cross-entropy with hard labels and KL divergence from teacher soft logits at temperature $\tau$:
\begin{equation}
  \mathcal{L}_{\text{KD}} = \alpha\,\mathcal{L}_{\text{CE}}(P_S, y_{\text{hard}}) + (1-\alpha)\,\tau^2\,\mathbb{KL}\!\left[P_T^{(\tau)} \| P_S^{(\tau)}\right]
\end{equation}
where $P^{(\tau)}_k = \text{softmax}(z_k/\tau)$. Temperature $\tau > 1$ softens distributions to expose similarity structure between classes not visible in hard labels.

\subsection{Architecture Flow Diagram}

\begin{figure}[h]
\centering
\begin{tikzpicture}[node distance=0.7cm and 1.0cm]
  \node[datblk] (x) {Input $x$};
  \node[frzblk, above right=0.5cm and 1.2cm of x] (teacher) {Teacher $P_T$\\(frozen, large)};
  \node[trnblk, below right=0.5cm and 1.2cm of x] (student) {Student $P_S$\\(trainable, small)};
  \node[lossblk, right=3.0cm of x] (kl) {KL Div\\$\tau^2\,\text{KL}[P_T^\tau \| P_S^\tau]$};
  \node[lossblk, below=0.7cm of kl] (ce) {CE Loss\\$\mathcal{L}_{\text{CE}}(P_S, y)$};
  \node[mrgblk, right=2.0cm of $(kl)!0.5!(ce)$] (total) {$\mathcal{L}_{\text{KD}}$\\(combined)};

  \draw[arr] (x.east) |- (teacher.west);
  \draw[arr] (x.east) |- (student.west);
  \draw[arr] (teacher.east) -| (kl.north);
  \draw[arr] (student.east) -| (kl.south);
  \draw[arr] (student.east) -| (ce.west);
  \draw[arr] (kl.east) -| (total.north);
  \draw[arr] (ce.east) -| (total.south);
\end{tikzpicture}
\caption{Knowledge Distillation: student learns from both hard labels (CE) and teacher soft logits (KL) simultaneously.}
\end{figure}

\subsection{State Transition Table}
\begin{table}[h]\centering
\begin{tabular}{lll}\toprule
\textbf{Property} & \textbf{Before} & \textbf{After}\\\midrule
Student accuracy & Low (random) & Near-teacher accuracy\\
Training signal & Hard labels only & Hard labels + soft teacher logits\\
Temperature $\tau$ & N/A & $>1$ (typically 2--5)\\
Dark knowledge transfer & None & Via inter-class similarities\\
Student size & Large (teacher) & Small (${\sim}10$--$100\times$ fewer params)\\\bottomrule
\end{tabular}
\caption{State properties before and after Knowledge Distillation.}
\end{table}

\newpage
%% ------------------------------------------------------------
\section{Response Distillation}

\subsection{Mathematical Foundation}
Response Distillation trains the student by directly imitating the teacher's \emph{full output sequence} rather than soft logits. The student minimises cross-entropy against teacher-generated responses:
\begin{equation}
  \mathcal{L}_{\text{RD}} = -\sum_{(x,y_T) \in \mathcal{D}_T} \sum_t \log P_S(y_{T,t} \mid x, y_{T,<t})
\end{equation}
where $y_T \sim P_T(\cdot \mid x)$ are teacher completions. This is simpler than logit matching (no access to teacher logits required) and enables distillation from black-box APIs.

\subsection{Architecture Flow Diagram}

\begin{figure}[h]
\centering
\begin{tikzpicture}[node distance=0.7cm and 1.1cm]
  \node[datblk] (prompts) {Prompts $\{x_i\}$};
  \node[frzblk, right=of prompts] (teacher) {Teacher $P_T$\\(API or local)};
  \node[datblk, right=of teacher] (responses) {Teacher\\Responses $y_T$};
  \node[trnblk, right=of responses] (student) {Student $P_S$\\SFT on $y_T$};
  \node[lossblk, below=0.9cm of student] (loss) {CE Loss\\$-\log P_S(y_T|x)$};
  \node[mrgblk, left=1.4cm of loss] (out) {Distilled\\Student};

  \draw[arr] (prompts)--(teacher);
  \draw[arr] (teacher)--(responses);
  \draw[arr] (responses)--(student);
  \draw[arr] (student)--(loss);
  \draw[arr] (loss)--(out);

  \node[font=\tiny,red,align=center,below=0.2cm of teacher] {black-box OK\\(only outputs needed)};
\end{tikzpicture}
\caption{Response Distillation: teacher generates responses; student trained via SFT to imitate those outputs.}
\end{figure}

\subsection{State Transition Table}
\begin{table}[h]\centering
\begin{tabular}{lll}\toprule
\textbf{Property} & \textbf{Before (KD)} & \textbf{After (Response KD)}\\\midrule
Teacher access required & Logits (white-box) & Outputs only (black-box OK)\\
Loss type & KL on logits & CE on token sequence\\
Sequence-level knowledge & Partial & Full (complete generation)\\
Implementation complexity & High & Same as SFT\\
Applicable to proprietary APIs & No & Yes\\\bottomrule
\end{tabular}
\caption{State properties comparing logit-level KD and Response Distillation.}
\end{table}

\newpage
%% ------------------------------------------------------------
\section{Chain-of-Thought Distillation}

\subsection{Mathematical Foundation}
CoT Distillation transfers the teacher's \emph{reasoning traces} to the student, so the student learns not just the answer but how to think step by step. The student is trained on $(x, c_T, y_T)$ triplets where $c_T$ is the teacher-generated chain:
\begin{equation}
  \mathcal{L}_{\text{CoTD}} = -\sum_{(x,c_T,y_T)} \left[\sum_m \log P_S(c_{T,m} \mid x, c_{T,<m}) + \sum_t \log P_S(y_{T,t} \mid x, c_T, y_{T,<t})\right]
\end{equation}
This allows a small student to match the accuracy of a large teacher that reasons step by step.

\subsection{Architecture Flow Diagram}

\begin{figure}[h]
\centering
\begin{tikzpicture}[node distance=0.7cm and 1.0cm]
  \node[datblk] (x) {Problem $x$};
  \node[frzblk, right=of x] (teacher) {Teacher $P_T$};
  \node[grayblk, right=of teacher] (cot) {Chain $c_T$\\+ Answer $y_T$};
  \node[trnblk, below=1.0cm of teacher] (student) {Student $P_S$};
  \node[lossblk, right=of student] (loss) {CoT + Answer\\CE Loss};
  \node[mrgblk, right=of loss] (out) {Reasoning\\Student};

  \draw[arr] (x)--(teacher);
  \draw[arr] (teacher)--(cot);
  \draw[arr] (cot.south) -- ++(0,-0.3) -| (loss.north);
  \draw[arr] (x.south) -- ++(0,-0.7) -| (student.west);
  \draw[arr] (student)--(loss);
  \draw[arr] (loss)--(out);
\end{tikzpicture}
\caption{CoT Distillation: teacher's reasoning chains and answers jointly train student to produce similar traces.}
\end{figure}

\subsection{State Transition Table}
\begin{table}[h]\centering
\begin{tabular}{lll}\toprule
\textbf{Property} & \textbf{Before} & \textbf{After}\\\midrule
Student reasoning & Direct answer & Step-by-step chain\\
Training target & Answer only & Chain + answer\\
Reasoning data source & None & Teacher-generated\\
Multi-step problem accuracy & Low & Significantly higher\\
Output verbosity & Terse & Explanatory\\\bottomrule
\end{tabular}
\caption{State properties before and after CoT Distillation.}
\end{table}

\newpage
%% ------------------------------------------------------------
\section{Policy Distillation}

\subsection{Mathematical Foundation}
Policy Distillation minimises the KL divergence between a teacher policy $\pi_T$ and student policy $\pi_S$ over the full output distribution at every decoding step, not just the final token:
\begin{equation}
  \mathcal{L}_{\text{PD}} = \mathbb{E}_{x \sim \mathcal{D}} \sum_t \mathbb{KL}\!\left[\pi_T(\cdot \mid x, y_{<t}) \| \pi_S(\cdot \mid x, y_{<t})\right]
\end{equation}
Unlike response distillation (which uses sampled strings), policy distillation uses teacher token-level distributions and requires teacher logit access. This aligns the student's generation \emph{process} with the teacher's, not just its outputs.

\subsection{Architecture Flow Diagram}

\begin{figure}[h]
\centering
\begin{tikzpicture}[node distance=0.7cm and 1.0cm]
  \node[datblk] (x) {Context $(x, y_{<t})$};
  \node[frzblk, above right=0.5cm and 1.3cm of x] (tpol) {Teacher\\$\pi_T(\cdot|x,y_{<t})$\\full distribution};
  \node[trnblk, below right=0.5cm and 1.3cm of x] (spol) {Student\\$\pi_S(\cdot|x,y_{<t})$\\full distribution};
  \node[lossblk, right=3.2cm of x] (kl) {KL$[\pi_T \| \pi_S]$\\at every step $t$};
  \node[mrgblk, right=of kl] (out) {Aligned\\Policy $\pi_S^*$};

  \draw[arr] (x.east) |- (tpol.west);
  \draw[arr] (x.east) |- (spol.west);
  \draw[arr] (tpol.east) -| (kl.north);
  \draw[arr] (spol.east) -| (kl.south);
  \draw[arr] (kl)--(out);

  \node[red,font=\tiny,align=center,below=0.2cm of tpol] {logits required\\(white-box)};
\end{tikzpicture}
\caption{Policy Distillation: KL divergence between teacher and student distributions minimised at every token step.}
\end{figure}

\subsection{State Transition Table}
\begin{table}[h]\centering
\begin{tabular}{lll}\toprule
\textbf{Property} & \textbf{Before (Response KD)} & \textbf{After (Policy KD)}\\\midrule
Teacher signal & Sampled string & Full token distribution\\
KL divergence & N/A & Minimised per step\\
Teacher access & Black-box & White-box (logits)\\
Process alignment & Output only & Step-by-step\\
Overconfidence risk & Higher & Lower (distribution match)\\\bottomrule
\end{tabular}
\caption{State properties comparing Response and Policy Distillation.}
\end{table}

\newpage
%% ------------------------------------------------------------
\section{Reward Distillation}

\subsection{Mathematical Foundation}
Reward Distillation compresses a large teacher reward model $r_T$ into a smaller student reward model $r_S$ by regression on teacher-generated scores:
\begin{equation}
  \mathcal{L}_{\text{RewardD}} = \mathbb{E}_{(x,y)}\bigl[(r_T(x,y) - r_S(x,y))^2\bigr]
\end{equation}
The student can also be trained with a ranking-consistency loss to preserve relative orderings even if absolute score values differ:
\begin{equation}
  \mathcal{L}_{\text{rank}} = -\log\sigma\!\left(r_S(x,y_w) - r_S(x,y_l)\right) \text{ whenever } r_T(x,y_w) > r_T(x,y_l)
\end{equation}

\subsection{Architecture Flow Diagram}

\begin{figure}[h]
\centering
\begin{tikzpicture}[node distance=0.7cm and 1.0cm]
  \node[datblk] (pairs) {$(x, y)$ Pairs};
  \node[frzblk, right=of pairs] (trm) {Teacher RM $r_T$\\(frozen, large)};
  \node[grayblk, right=of trm] (scores) {Teacher\\Scores $r_T(x,y)$};
  \node[trnblk, below=1.0cm of trm] (srm) {Student RM $r_S$\\(small, trainable)};
  \node[lossblk, right=of srm] (loss) {MSE +\\Ranking Loss};
  \node[mrgblk, right=of loss] (out) {Compact\\RM $r_S^*$};

  \draw[arr] (pairs)--(trm);
  \draw[arr] (trm)--(scores);
  \draw[arr] (scores.south) -- ++(0,-0.3) -| (loss.north);
  \draw[arr] (pairs.south) -- ++(0,-0.7) -| (srm.west);
  \draw[arr] (srm)--(loss);
  \draw[arr] (loss)--(out);
\end{tikzpicture}
\caption{Reward Distillation: small student reward model trained via regression on large teacher reward scores.}
\end{figure}

\subsection{State Transition Table}
\begin{table}[h]\centering
\begin{tabular}{lll}\toprule
\textbf{Property} & \textbf{Before} & \textbf{After}\\\midrule
Reward model size & Large teacher & Small student\\
Inference latency & Slow & Fast\\
Score accuracy & High (teacher) & Near-teacher (distilled)\\
Training objective & Human preferences & Teacher score regression\\
Deployment cost & High & Low\\\bottomrule
\end{tabular}
\caption{State properties before and after Reward Distillation.}
\end{table}

\newpage
%% ------------------------------------------------------------
\section{Retrieval Distillation}

\subsection{Mathematical Foundation}
Retrieval Distillation trains a student retriever $q_S$ to match the relevance judgments of a larger teacher retriever $q_T$ (or cross-encoder) without requiring access to teacher gradients. The student minimises the KL divergence between teacher and student retrieval distributions over a document set $\mathcal{C}$:
\begin{equation}
  \mathcal{L}_{\text{RD}} = \mathbb{KL}\!\left[P_T(\cdot \mid q) \| P_S(\cdot \mid q)\right],\quad P(d \mid q) \propto \exp(q_\phi(q)^\top e(d))
\end{equation}
where $q_\phi$ is the query encoder and $e(d)$ is the document embedding.

\subsection{Architecture Flow Diagram}

\begin{figure}[h]
\centering
\begin{tikzpicture}[node distance=0.7cm and 1.1cm]
  \node[datblk] (q) {Query $q$};
  \node[frzblk, above right=0.5cm and 1.2cm of q] (tret) {Teacher\\Retriever $q_T$\\(cross-encoder)};
  \node[trnblk, below right=0.5cm and 1.2cm of q] (sret) {Student\\Retriever $q_S$\\(bi-encoder)};
  \node[datblk, right=3.2cm of q] (corpus) {Document\\Corpus $\mathcal{C}$};
  \node[lossblk, right=4.5cm of q] (kl) {KL[$P_T \| P_S$]\\over doc ranks};
  \node[mrgblk, right=of kl] (out) {Efficient\\Retriever};

  \draw[arr] (q.east) |- (tret.west);
  \draw[arr] (q.east) |- (sret.west);
  \draw[arr] (tret.east) -| (kl.north);
  \draw[arr] (sret.east) -| (kl.south);
  \draw[arr] (corpus.west) -- (kl.east|-corpus) -- (kl.east);
  \draw[arr] (kl)--(out);

  \node[red,font=\tiny,align=center,below=0.2cm of sret] {fast at inference\\(ANN search)};
\end{tikzpicture}
\caption{Retrieval Distillation: student bi-encoder trained to match teacher cross-encoder relevance distributions.}
\end{figure}

\subsection{State Transition Table}
\begin{table}[h]\centering
\begin{tabular}{lll}\toprule
\textbf{Property} & \textbf{Before (Teacher)} & \textbf{After (Student)}\\\midrule
Retriever type & Cross-encoder (slow) & Bi-encoder (fast ANN)\\
Query latency & $O(|\mathcal{C}|)$ & $O(\log|\mathcal{C}|)$\\
Ranking quality & Very high & Near-teacher\\
Index requirement & None & Pre-computed embeddings\\
Training signal & Human judgments & Teacher soft scores\\\bottomrule
\end{tabular}
\caption{State properties comparing teacher cross-encoder and student bi-encoder after retrieval distillation.}
\end{table}

\newpage
%% ============================================================
\part{Model Compression}
%% ============================================================

\section{Quantization-Aware Training (QAT)}

\subsection{Mathematical Foundation}
QAT simulates quantization during training using a straight-through estimator (STE) for gradients. For weight $w$, the quantized value at $b$ bits with scale $s$:
\begin{equation}
  \hat{w} = \text{clamp}\!\left(\text{round}(w/s), -2^{b-1}, 2^{b-1}-1\right) \cdot s
\end{equation}
During the backward pass the STE treats quantization as identity for gradient flow:
\begin{equation}
  \frac{\partial \mathcal{L}}{\partial w} \approx \frac{\partial \mathcal{L}}{\partial \hat{w}} \cdot \mathbf{1}\!\left[|w| \leq s \cdot 2^{b-1}\right]
\end{equation}
This makes the model aware of quantization error during training, yielding better accuracy than post-training quantization.

\subsection{Architecture Flow Diagram}

\begin{figure}[h]
\centering
\begin{tikzpicture}[node distance=0.7cm and 1.1cm]
  \node[datblk] (x) {Input $x$};
  \node[trnblk, right=of x] (fp32) {FP32 Weights\\(master copy)};
  \node[grayblk, right=of fp32] (fakeq) {Fake Quantize\\$w \to \hat{w}$};
  \node[frzblk, right=of fakeq] (fwd) {Forward Pass\\(with $\hat{w}$)};
  \node[lossblk, right=of fwd] (loss) {Task Loss\\$\mathcal{L}$};
  \node[trnblk, below=0.9cm of loss] (ste) {STE Backward\\$\partial\mathcal{L}/\partial w$};
  \node[mrgblk, left=3.6cm of ste] (deploy) {Deploy\\INT$b$ Model};

  \draw[arr] (x)--(fp32);
  \draw[arr] (fp32)--(fakeq);
  \draw[arr] (fakeq)--(fwd);
  \draw[arr] (fwd)--(loss);
  \draw[arr] (loss)--(ste);
  \draw[arr] (ste) -- node[below,font=\tiny]{update FP32 master} (deploy);
  \draw[arr,dashed,green!60!black] (ste.west) -- (fp32.south|-ste) -- (fp32.south);
\end{tikzpicture}
\caption{QAT: fake quantization in forward pass; STE allows gradients to flow to FP32 master weights.}
\end{figure}

\subsection{State Transition Table}
\begin{table}[h]\centering
\begin{tabular}{lll}\toprule
\textbf{Property} & \textbf{Before (PTQ)} & \textbf{After (QAT)}\\\midrule
Quantization awareness & None during training & Simulated in forward pass\\
Accuracy loss & 1--5\% degradation & $<$1\% typical\\
Calibration data & Small set required & Full training data\\
Weight format & FP32 & FP32 master, INT$b$ deploy\\
Gradient flow & Normal & Via STE\\\bottomrule
\end{tabular}
\caption{State properties comparing Post-Training Quantization and QAT.}
\end{table}

\newpage
%% ------------------------------------------------------------
\section{Pruning}

\subsection{Mathematical Foundation}
Pruning removes weights whose contribution to the model output is below a threshold, producing a sparse network. Magnitude-based unstructured pruning removes the fraction $p$ of weights with the smallest absolute values:
\begin{equation}
  m_i = \mathbf{1}[|w_i| \geq \tau_p],\quad W_{\text{sparse}} = W \odot m
\end{equation}
where $\tau_p$ is chosen to achieve sparsity ratio $s = 1 - \|m\|_1 / \|m\|_0$. Structured pruning removes entire heads or neurons (rows/columns) and is hardware-friendly:
\begin{equation}
  \text{importance}(h) = \sum_{i \in h} |w_i| \cdot |\nabla_{w_i}\mathcal{L}|
\end{equation}

\subsection{Architecture Flow Diagram}

\begin{figure}[h]
\centering
\begin{tikzpicture}[node distance=0.7cm and 1.0cm]
  \node[frzblk] (dense) {Dense Model\\$W$ (all weights)};
  \node[grayblk, right=of dense] (score) {Importance\\Scoring};
  \node[orangeblk, right=of score] (mask) {Binary\\Mask $m$};
  \node[trnblk, right=of mask] (sparse) {Sparse Model\\$W \odot m$};
  \node[lossblk, below=0.9cm of sparse] (finetune) {Fine-tune to\\Recover Acc.};
  \node[mrgblk, left=2.2cm of finetune] (out) {Pruned\\Model};

  \draw[arr] (dense)--(score);
  \draw[arr] (score)--(mask);
  \draw[arr] (mask)--(sparse);
  \draw[arr] (sparse)--(finetune);
  \draw[arr] (finetune)--(out);

  \node[font=\tiny,below=0.2cm of mask] {sparsity $s = 50$--$90$\%};
\end{tikzpicture}
\caption{Pruning: importance scores determine binary mask; sparse model fine-tuned to recover accuracy.}
\end{figure}

\subsection{State Transition Table}
\begin{table}[h]\centering
\begin{tabular}{lll}\toprule
\textbf{Property} & \textbf{Before} & \textbf{After}\\\midrule
Model sparsity & 0\% & 50--90\%\\
Parameter count & $|\theta|$ & $(1-s)|\theta|$\\
Inference FLOPs & Baseline & Reduced (structured)\\
Accuracy & Baseline & Slightly reduced\\
Memory footprint & Full & Reduced by $s$\\\bottomrule
\end{tabular}
\caption{State properties before and after pruning.}
\end{table}

\newpage
%% ------------------------------------------------------------
\section{Post-Training Quantization (PTQ)}

\subsection{Mathematical Foundation}
PTQ converts a trained FP32 model to reduced precision (INT8, INT4) without additional training. Per-channel scale factors $s_c$ are calibrated on a small unlabelled dataset $\mathcal{D}_{\text{cal}}$:
\begin{equation}
  s_c = \frac{\max_{x \in \mathcal{D}_{\text{cal}}}|a_c(x)|}{2^{b-1}-1},\quad \hat{a}_c = \text{round}(a_c / s_c) \cdot s_c
\end{equation}
GPTQ computes optimal quantization for each weight row by minimising the second-order reconstruction error using the Hessian:
\begin{equation}
  \min_{\hat{W}} \|\hat{W} X - W X\|_F^2 \approx \sum_i (w_i - \hat{w}_i)^2 [H^{-1}]_{ii}
\end{equation}

\subsection{Architecture Flow Diagram}

\begin{figure}[h]
\centering
\begin{tikzpicture}[node distance=0.7cm and 1.0cm]
  \node[frzblk] (fp32) {FP32 Model\\(pretrained)};
  \node[datblk, right=of fp32] (cal) {Calibration\\Data $\mathcal{D}_{\text{cal}}$};
  \node[grayblk, right=of cal] (hess) {Hessian /\\Scale Est.};
  \node[orangeblk, right=of hess] (quant) {Round Weights\\to INT$b$};
  \node[mrgblk, right=of quant] (int) {INT$b$ Model\\(fast inference)};

  \draw[arr] (fp32)--(cal);
  \draw[arr] (cal)--(hess);
  \draw[arr] (hess)--(quant);
  \draw[arr] (quant)--(int);
  \draw[arr] (fp32.south) -- ++(0,-0.5) -| (quant.south);

  \node[font=\tiny,below=0.2cm of cal] {small unlabelled};
\end{tikzpicture}
\caption{PTQ: calibration statistics guide per-channel scales; weights rounded to INT$b$ without retraining.}
\end{figure}

\subsection{State Transition Table}
\begin{table}[h]\centering
\begin{tabular}{lll}\toprule
\textbf{Property} & \textbf{Before} & \textbf{After}\\\midrule
Weight precision & FP32 / BF16 & INT8 / INT4\\
Memory usage & $4 \times n$ bytes & $1$--$0.5 \times n$ bytes\\
Retraining required & N/A & None (calibration only)\\
Accuracy loss & 0\% & 0.5--3\% (INT8), higher INT4\\
Calibration data & Not needed & Small sample ($128$--$512$ examples)\\\bottomrule
\end{tabular}
\caption{State properties before and after Post-Training Quantization.}
\end{table}

\newpage
%% ============================================================
\part{Agentic Training}
%% ============================================================

\section{Tool-Use Training}

\subsection{Mathematical Foundation}
Tool-Use Training teaches the model to generate structured tool invocations interleaved with reasoning (ReAct-style). The training data consists of trajectories $\tau = (t_1, a_1, o_1, \ldots, t_N, a_N, o_N, y)$ where $t_i$ is a thought, $a_i$ is a tool call, and $o_i$ is the observation. The model is trained to predict thought and action tokens conditioned on prior context and observations:
\begin{equation}
  \mathcal{L}_{\text{tool}} = -\sum_{i=1}^N \sum_j \log P_\theta(t_{i,j} \mid \text{context}_{<i}) + \log P_\theta(a_{i,j} \mid \text{context}_{<i}, t_i)
\end{equation}

\subsection{Architecture Flow Diagram}

\begin{figure}[h]
\centering
\begin{tikzpicture}[node distance=0.5cm and 0.8cm]
  \node[datblk,minimum width=1.6cm] (q) {Task $q$};
  \node[trnblk,minimum width=1.6cm, right=0.7cm of q] (t1) {Thought\\$t_1$};
  \node[orangeblk,minimum width=1.6cm, right=0.6cm of t1] (a1) {Tool Call\\$a_1$};
  \node[cyanblk,minimum width=1.6cm, right=0.6cm of a1] (o1) {Obs.\\$o_1$};
  \node[trnblk,minimum width=1.6cm, right=0.6cm of o1] (t2) {Thought\\$t_2$};
  \node[orangeblk,minimum width=1.6cm, right=0.6cm of t2] (a2) {Tool Call\\$a_2$};
  \node[mrgblk,minimum width=1.6cm, right=0.6cm of a2] (ans) {Answer\\$y$};

  \foreach \a/\b in {q/t1, t1/a1, a1/o1, o1/t2, t2/a2, a2/ans}
    \draw[arr] (\a)--(\b);

  \node[red,font=\tiny,below=0.3cm of a1] {trained};
  \node[red,font=\tiny,below=0.3cm of t1] {trained};
  \node[blue!60!black,font=\tiny,below=0.3cm of o1] {env reply};
\end{tikzpicture}
\caption{Tool-Use Training: thought-action-observation trajectory; model trained on thought and action tokens.}
\end{figure}

\subsection{State Transition Table}
\begin{table}[h]\centering
\begin{tabular}{lll}\toprule
\textbf{Property} & \textbf{Before} & \textbf{After}\\\midrule
Output format & Free text & Thought + tool call interleaved\\
Tool invocation & None & Structured API calls\\
Observation integration & N/A & Conditioned on tool results\\
Multi-step planning & Weak & Strong\\
Training data type & Static QA & Trajectory (thought, action, obs)\\\bottomrule
\end{tabular}
\caption{State properties before and after Tool-Use Training.}
\end{table}

\newpage
%% ------------------------------------------------------------
\section{Function Calling Training}

\subsection{Mathematical Foundation}
Function Calling Training specifically trains the model to produce well-formed JSON function-call objects given a natural-language request and a tool schema. The model maps input $x$ and schema $S = \{f_1,\ldots,f_K\}$ to a structured call $c = \{$\texttt{name}: $f_k$, \texttt{args}: $\{a_1{:}v_1,\ldots\}\}$:
\begin{equation}
  \mathcal{L}_{\text{FC}} = -\log P_\theta(c \mid x, S)
\end{equation}
where $c$ is the ground-truth function call. Schema-aware token masking ensures the model only generates valid JSON that matches the function signature, reducing hallucination of non-existent arguments.

\subsection{Architecture Flow Diagram}

\begin{figure}[h]
\centering
\begin{tikzpicture}[node distance=0.7cm and 1.1cm]
  \node[datblk] (x) {User Request $x$};
  \node[grayblk, above=0.7cm of x] (schema) {Tool Schema $S$\\(JSON spec)};
  \node[frzblk, right=1.6cm of $(x)!0.5!(schema)$] (model) {LLM\\$P_\theta$};
  \node[trnblk, right=of model] (json) {JSON Call\\$c=\{$name, args$\}$};
  \node[lossblk, below=0.9cm of json] (val) {Schema\\Validation};
  \node[mrgblk, left=1.4cm of val] (exec) {Execute\\Function};

  \draw[arr] (x.east) -| (model.south);
  \draw[arr] (schema.east) -| (model.north);
  \draw[arr] (model)--(json);
  \draw[arr] (json)--(val);
  \draw[arr] (val)--(exec);
\end{tikzpicture}
\caption{Function Calling Training: model receives request + schema, outputs structured JSON call validated against spec.}
\end{figure}

\subsection{State Transition Table}
\begin{table}[h]\centering
\begin{tabular}{lll}\toprule
\textbf{Property} & \textbf{Before} & \textbf{After}\\\midrule
Output format & Natural language & Structured JSON\\
Schema adherence & None & Enforced\\
Argument hallucination & High & Low\\
API integration & Manual parsing & Direct execution\\
Training data & Text completions & (request, schema, call) triplets\\\bottomrule
\end{tabular}
\caption{State properties before and after Function Calling Training.}
\end{table}

\newpage
%% ------------------------------------------------------------
\section{Multi-Agent Training}

\subsection{Mathematical Foundation}
Multi-Agent Training trains a collection of specialised agents $\{\pi_1,\ldots,\pi_K\}$ and an orchestrator $\pi_0$ to collaboratively solve complex tasks. The orchestrator decomposes a task $T$ into sub-tasks $\{t_1,\ldots,t_K\}$ and assigns them to specialists:
\begin{equation}
  \pi_0: T \to \{(t_k, k)\}_{k=1}^K, \quad \pi_k: t_k \to y_k, \quad y = \pi_0(\{y_k\})
\end{equation}
Training uses task-completion reward at the full-task level, distributed to agents via credit assignment:
\begin{equation}
  r_k = r_{\text{task}} \cdot \mathbf{1}[\pi_k \text{ contributed to solution}]
\end{equation}

\subsection{Architecture Flow Diagram}

\begin{figure}[h]
\centering
\begin{tikzpicture}[node distance=0.6cm and 1.0cm]
  \node[datblk] (task) {Complex\\Task $T$};
  \node[frzblk,minimum height=1.4cm] (orch) at (3.2,0) {Orchestrator\\$\pi_0$};
  \node[trnblk,minimum width=1.6cm] (s1) at (6.2,1.2) {Specialist $\pi_1$\\(code)};
  \node[trnblk,minimum width=1.6cm] (s2) at (6.2,0.0) {Specialist $\pi_2$\\(search)};
  \node[trnblk,minimum width=1.6cm] (s3) at (6.2,-1.2) {Specialist $\pi_K$\\(math)};
  \node[mrgblk] (agg) at (9.2,0) {Aggregate\\Answer $y$};

  \draw[arr] (task)--(orch);
  \draw[arr] (orch.east) |- (s1.west);
  \draw[arr] (orch.east) -- (s2.west);
  \draw[arr] (orch.east) |- (s3.west);
  \draw[arr] (s1.east) -| (agg.north);
  \draw[arr] (s2)--(agg);
  \draw[arr] (s3.east) -| (agg.south);
\end{tikzpicture}
\caption{Multi-Agent Training: orchestrator decomposes task and assigns sub-tasks to specialised agents.}
\end{figure}

\subsection{State Transition Table}
\begin{table}[h]\centering
\begin{tabular}{lll}\toprule
\textbf{Property} & \textbf{Before} & \textbf{After}\\\midrule
Task handling & Single monolithic model & Orchestrator + specialists\\
Specialisation & None & Per-domain agents\\
Credit assignment & N/A & Per-agent contribution\\
Task complexity ceiling & Single-context limit & Decomposable to multiple contexts\\
Coordination & N/A & Learned delegation protocol\\\bottomrule
\end{tabular}
\caption{State properties before and after Multi-Agent Training.}
\end{table}

\newpage
%% ============================================================
\part{Multimodal Training}
%% ============================================================

\section{Vision-Language Training}

\subsection{Mathematical Foundation}
Vision-Language Training jointly trains a visual encoder $f_v$ and language model $f_l$ to process image-text pairs. A projection layer $P$ aligns visual features to the LLM's token embedding space:
\begin{equation}
  v_{\text{tokens}} = P(f_v(I)),\quad \tilde{x} = [v_{\text{tokens}};\, E(x_{\text{text}})]
\end{equation}
Training uses a masked language modelling or auto-regressive loss on the text tokens conditioned on visual tokens:
\begin{equation}
  \mathcal{L}_{\text{VL}} = -\sum_t \log P_{f_l}(x_t \mid v_{\text{tokens}}, x_{<t})
\end{equation}

\subsection{Architecture Flow Diagram}

\begin{figure}[h]
\centering
\begin{tikzpicture}[node distance=0.7cm and 1.1cm]
  \node[datblk] (img) {Image $I$};
  \node[datblk, below=0.8cm of img] (txt) {Text $x$};
  \node[frzblk, right=of img] (venc) {Visual\\Encoder $f_v$\\(ViT/CLIP)};
  \node[grayblk, right=of venc] (proj) {Projection\\Layer $P$};
  \node[trnblk, right=of txt] (emb) {Text\\Embeddings};
  \node[frzblk, right=3.5cm of $(proj)!0.5!(emb)$] (llm) {LLM\\$f_l$};
  \node[lossblk, right=of llm] (loss) {LM Loss\\on text};

  \draw[arr] (img)--(venc);
  \draw[arr] (venc)--(proj);
  \draw[arr] (txt)--(emb);
  \draw[arr] (proj.east) -| (llm.north);
  \draw[arr] (emb.east) -| (llm.south);
  \draw[arr] (llm)--(loss);
\end{tikzpicture}
\caption{Vision-Language Training: visual encoder outputs projected to LLM token space; joint auto-regressive loss.}
\end{figure}

\subsection{State Transition Table}
\begin{table}[h]\centering
\begin{tabular}{lll}\toprule
\textbf{Property} & \textbf{Before} & \textbf{After}\\\midrule
Input modality & Text only & Text + images\\
Visual understanding & None & Strong (via encoder)\\
Projection layer & Not present & Trained to align modalities\\
Alignment quality & N/A & Measured by VQA, captioning\\
Training data & Text corpora & Image-caption / VQA pairs\\\bottomrule
\end{tabular}
\caption{State properties before and after Vision-Language Training.}
\end{table}

\newpage
%% ------------------------------------------------------------
\section{Grounding Training}

\subsection{Mathematical Foundation}
Grounding Training teaches the model to link text spans to image regions (bounding boxes). Given a referring expression $e$ and image $I$, the model predicts box coordinates $(x_1,y_1,x_2,y_2)$ normalised to $[0,1]^4$:
\begin{equation}
  \hat{b} = f_\theta(I, e),\quad \mathcal{L}_{\text{GND}} = \text{L1}(\hat{b}, b^*) + \lambda(1 - \text{IoU}(\hat{b}, b^*))
\end{equation}
where $b^*$ is the ground-truth box and IoU is intersection-over-union. The model generates boxes as text tokens (serialised coordinates), enabling unified seq2seq training.

\subsection{Architecture Flow Diagram}

\begin{figure}[h]
\centering
\begin{tikzpicture}[node distance=0.7cm and 1.1cm]
  \node[datblk] (img) {Image $I$};
  \node[datblk, below=0.8cm of img] (expr) {Referring\\Expr.\ $e$};
  \node[frzblk, right=1.6cm of $(img)!0.5!(expr)$] (vlm) {VLM\\$f_\theta$};
  \node[trnblk, right=of vlm] (box) {Predicted Box\\$(x_1,y_1,x_2,y_2)$};
  \node[lossblk, right=of box] (loss) {L1 +\\IoU Loss};
  \node[datblk, above=0.7cm of loss] (gt) {GT Box\\$b^*$};

  \draw[arr] (img.east) -| (vlm.north);
  \draw[arr] (expr.east) -| (vlm.south);
  \draw[arr] (vlm)--(box);
  \draw[arr] (box)--(loss);
  \draw[arr] (gt.south)--(loss.north);
\end{tikzpicture}
\caption{Grounding Training: VLM predicts bounding box coordinates from image and referring expression.}
\end{figure}

\subsection{State Transition Table}
\begin{table}[h]\centering
\begin{tabular}{lll}\toprule
\textbf{Property} & \textbf{Before} & \textbf{After}\\\midrule
Output type & Text description & Bounding box coordinates\\
Spatial understanding & Implicit & Explicit (pixel-level)\\
Loss type & Cross-entropy & L1 + IoU\\
Grounding ability & None & Text-to-region linking\\
Training data & Captions & (image, expression, bbox) triplets\\\bottomrule
\end{tabular}
\caption{State properties before and after Grounding Training.}
\end{table}

\newpage
%% ============================================================
\part{Knowledge Injection}
%% ============================================================

\section{Retrieval-Augmented Training}

\subsection{Mathematical Foundation}
Retrieval-Augmented Training (RAT) trains the model to use retrieved documents $\mathcal{R}(x) = \{d_1,\ldots,d_K\}$ as context. The model generates $y$ conditioned on the query and retrieved set:
\begin{equation}
  P_\theta(y \mid x, \mathcal{R}(x)) = \prod_t P_\theta(y_t \mid x, d_1,\ldots,d_K, y_{<t})
\end{equation}
Training includes both retrieval-present and retrieval-absent examples to prevent over-reliance on context. A closed-book loss component ensures the model retains parametric knowledge:
\begin{equation}
  \mathcal{L}_{\text{RAT}} = \lambda\,\mathcal{L}_{\text{open-book}} + (1-\lambda)\,\mathcal{L}_{\text{closed-book}}
\end{equation}

\subsection{Architecture Flow Diagram}

\begin{figure}[h]
\centering
\begin{tikzpicture}[node distance=0.7cm and 1.1cm]
  \node[datblk] (q) {Query $x$};
  \node[frzblk, above right=0.5cm and 1.2cm of q] (ret) {Retriever\\$\mathcal{R}(x)$};
  \node[datblk, right=1.2cm of q] (docs) {Docs\\$d_1..d_K$};
  \node[trnblk, right=of docs] (llm) {LLM\\$P_\theta$};
  \node[lossblk, right=of llm] (loss) {Open-book\\Loss};
  \node[lossblk, below=0.8cm of llm] (cbloss) {Closed-book\\Loss};
  \node[mrgblk, right=1.6cm of $(loss)!0.5!(cbloss)$] (out) {RAG-tuned\\Model};

  \draw[arr] (q.east) |- (ret.west);
  \draw[arr] (ret.east) -| (docs.north);
  \draw[arr] (q)--(docs);
  \draw[arr] (docs)--(llm);
  \draw[arr] (llm)--(loss);
  \draw[arr] (q.south) -- ++(0,-0.5) -| (cbloss.west);
  \draw[arr] (cbloss.east)-|(out.south);
  \draw[arr] (loss.east)-|(out.north);
\end{tikzpicture}
\caption{RAT: model trained on both retrieval-augmented and closed-book losses to balance open/parametric knowledge.}
\end{figure}

\subsection{State Transition Table}
\begin{table}[h]\centering
\begin{tabular}{lll}\toprule
\textbf{Property} & \textbf{Before} & \textbf{After}\\\midrule
Context awareness & Parametric only & Retrieved + parametric\\
Knowledge staleness & Fixed (training cutoff) & Dynamic (fresh retrieval)\\
Hallucination rate & Higher & Lower (grounded in docs)\\
Retriever dependency & None & Required at inference\\
Parametric retention & Full & Preserved (closed-book loss)\\\bottomrule
\end{tabular}
\caption{State properties before and after Retrieval-Augmented Training.}
\end{table}

\newpage
%% ------------------------------------------------------------
\section{Memory Tuning}

\subsection{Mathematical Foundation}
Memory Tuning trains the model to store and retrieve episodic memories $\mathcal{M} = \{(k_i, v_i)\}$ in an external key-value store. The model learns to write memories via a write head $w_\theta$ and retrieve via dot-product attention:
\begin{equation}
  \text{retrieve}(q) = \sum_i \text{softmax}(q^\top k_i / \sqrt{d}) \cdot v_i,\quad q = f_\theta^{\text{query}}(x)
\end{equation}
Training uses a distractor task: memory-relevant queries should prefer correct memory entries, trained with a contrastive retrieval loss:
\begin{equation}
  \mathcal{L}_{\text{mem}} = -\log \frac{\exp(q^\top k^+ / \tau)}{\sum_j \exp(q^\top k_j / \tau)}
\end{equation}

\subsection{Architecture Flow Diagram}

\begin{figure}[h]
\centering
\begin{tikzpicture}[node distance=0.7cm and 1.1cm]
  \node[datblk] (ep) {Episode\\$(x_t, y_t)$};
  \node[trnblk, right=of ep] (write) {Write Head\\$w_\theta$};
  \node[grayblk, right=of write, minimum height=2.0cm] (mem) {Memory\\Store\\$\mathcal{M}$};
  \node[datblk, below=1.2cm of ep] (query) {Future\\Query $x_{t'}$};
  \node[trnblk, right=of query] (read) {Read Head\\(dot-product)};
  \node[mrgblk, right=of read] (ctx) {Retrieved\\Context $v$};
  \node[lossblk, right=of ctx] (loss) {Contrastive\\Memory Loss};

  \draw[arr] (ep)--(write);
  \draw[arr] (write)--(mem);
  \draw[arr] (query)--(read);
  \draw[arr] (mem.south) -| (read.north);
  \draw[arr] (read)--(ctx);
  \draw[arr] (ctx)--(loss);
\end{tikzpicture}
\caption{Memory Tuning: write head stores episodes; read head retrieves via attention; contrastive loss trains both.}
\end{figure}

\subsection{State Transition Table}
\begin{table}[h]\centering
\begin{tabular}{lll}\toprule
\textbf{Property} & \textbf{Before} & \textbf{After}\\\midrule
Memory type & Parametric only & Parametric + episodic KV store\\
Context window & Fixed $T$ tokens & Unbounded (external memory)\\
Personalisation & None & Per-user episodic recall\\
Read/write mechanism & None & Learned attention-based\\
Long-term recall & Poor & Strong\\\bottomrule
\end{tabular}
\caption{State properties before and after Memory Tuning.}
\end{table}

\newpage
%% ------------------------------------------------------------
\section{Domain Knowledge Injection}

\subsection{Mathematical Foundation}
Domain Knowledge Injection embeds structured domain knowledge (ontologies, knowledge graphs, fact databases) into the model. Entities $e$ and relations $r$ from a KG are embedded alongside token embeddings and injected via cross-attention:
\begin{equation}
  h'_t = h_t + \text{CrossAttn}(h_t,\, \{e_i, r_{ij}\}_{(i,j) \in \mathcal{G}})
\end{equation}
Training uses a knowledge grounding loss that verifies generated text against retrieved facts:
\begin{equation}
  \mathcal{L}_{\text{KI}} = \mathcal{L}_{\text{LM}} + \mu\,\mathcal{L}_{\text{fact-check}}
\end{equation}

\subsection{Architecture Flow Diagram}

\begin{figure}[h]
\centering
\begin{tikzpicture}[node distance=0.7cm and 1.1cm]
  \node[datblk] (kg) {Knowledge\\Graph $\mathcal{G}$};
  \node[grayblk, right=of kg] (emb) {Entity/Relation\\Embeddings};
  \node[datblk, below=1.0cm of kg] (text) {Text Tokens\\$h_t$};
  \node[trnblk, right=of emb] (cross) {Cross-Attention\\Injection};
  \node[frzblk, right=of cross] (llm) {LLM\\Layers};
  \node[lossblk, right=of llm] (loss) {LM +\\Fact-check Loss};
  \node[mrgblk, below=0.9cm of loss] (out) {Knowledge-\\Grounded Model};

  \draw[arr] (kg)--(emb);
  \draw[arr] (emb)--(cross);
  \draw[arr] (text.east) -| (cross.south);
  \draw[arr] (cross)--(llm);
  \draw[arr] (llm)--(loss);
  \draw[arr] (loss)--(out);
\end{tikzpicture}
\caption{Domain Knowledge Injection: KG entity/relation embeddings injected into LLM via cross-attention.}
\end{figure}

\subsection{State Transition Table}
\begin{table}[h]\centering
\begin{tabular}{lll}\toprule
\textbf{Property} & \textbf{Before} & \textbf{After}\\\midrule
Knowledge source & Implicit (pretraining) & Explicit KG entities\\
Fact accuracy & Moderate & Significantly higher\\
Knowledge update & Retrain required & KG update sufficient\\
Cross-attention overhead & None & One extra module per layer\\
Domain coverage & General & Domain-specific + general\\\bottomrule
\end{tabular}
\caption{State properties before and after Domain Knowledge Injection.}
\end{table}

\newpage
%% ============================================================
\part{Deployment and Training Efficiency}
%% ============================================================

\section{Mixed Precision Training}

\subsection{Mathematical Foundation}
Mixed Precision Training maintains FP32 master weights but computes forward and backward passes in FP16/BF16. A loss scaling factor $S$ prevents underflow of small gradients:
\begin{equation}
  \tilde{\mathcal{L}} = S \cdot \mathcal{L},\quad g_{\text{FP32}} = \frac{1}{S} g_{\text{FP16}}
\end{equation}
BF16 has the same exponent range as FP32 (8 bits) but fewer mantissa bits (7 vs 23), making it more numerically stable without loss scaling:
\begin{equation}
  \text{memory}(\text{FP16}) = 2N \text{ bytes vs } \text{memory}(\text{FP32}) = 4N \text{ bytes}
\end{equation}

\subsection{Architecture Flow Diagram}

\begin{figure}[h]
\centering
\begin{tikzpicture}[node distance=0.7cm and 1.1cm]
  \node[frzblk] (master) {FP32 Master\\Weights $W$};
  \node[grayblk, right=of master] (cast) {Cast to\\FP16/BF16};
  \node[trnblk, right=of cast] (fwd) {Forward +\\Backward (FP16)};
  \node[lossblk, right=of fwd] (scale) {Loss Scale\\$S \cdot \mathcal{L}$};
  \node[trnblk, below=0.9cm of scale] (unscale) {Unscale Grad\\$g/S$};
  \node[mrgblk, left=3.6cm of unscale] (upd) {Adam Update\\(FP32 master)};

  \draw[arr] (master)--(cast);
  \draw[arr] (cast)--(fwd);
  \draw[arr] (fwd)--(scale);
  \draw[arr] (scale)--(unscale);
  \draw[arr] (unscale) -- node[below,font=\tiny]{check overflow} (upd);
  \draw[arr,dashed,green!60!black] (upd.north) |- (master.south) node[midway,below,font=\tiny]{update};
\end{tikzpicture}
\caption{Mixed Precision Training: FP16 compute with FP32 master weights; loss scaling prevents gradient underflow.}
\end{figure}

\subsection{State Transition Table}
\begin{table}[h]\centering
\begin{tabular}{lll}\toprule
\textbf{Property} & \textbf{Before (FP32)} & \textbf{After (Mixed FP16)}\\\midrule
Compute precision & FP32 & FP16/BF16\\
Memory footprint & $4N$ bytes & $2N + 4N$ (half + master)\\
Training speed & Baseline & $2$--$3\times$ faster\\
Numerical stability & High & Managed via loss scaling\\
GPU tensor cores & Unused & Utilised (FP16/BF16 native)\\\bottomrule
\end{tabular}
\caption{State properties before and after Mixed Precision Training.}
\end{table}

\newpage
%% ------------------------------------------------------------
\section{Gradient Checkpointing}

\subsection{Mathematical Foundation}
Gradient Checkpointing trades compute for memory by not storing all intermediate activations during the forward pass. Instead, activations at checkpoint nodes $\mathcal{C} \subset \{1,\ldots,L\}$ are stored; all others are recomputed during the backward pass:
\begin{equation}
  \text{Memory}_{\text{GC}} = O(\sqrt{L} \cdot d),\quad \text{vs }\; O(L \cdot d) \text{ without GC}
\end{equation}
For $L$ layers, optimal checkpointing every $\sqrt{L}$ layers minimises memory while adding one extra forward pass overhead:
\begin{equation}
  \text{Compute}_{\text{GC}} \approx 2 \cdot \text{Compute}_{\text{standard}}
\end{equation}

\subsection{Architecture Flow Diagram}

\begin{figure}[h]
\centering
\begin{tikzpicture}[node distance=0.5cm and 0.8cm]
  \node[datblk,minimum width=1.4cm] (in) {Input};
  \node[trnblk,minimum width=1.4cm, right=0.6cm of in] (l1) {Layer 1\\{\tiny (stored)}};
  \node[grayblk,minimum width=1.4cm, right=0.5cm of l1] (l2) {Layer 2\\{\tiny (drop)}};
  \node[grayblk,minimum width=1.4cm, right=0.5cm of l2] (l3) {Layer 3\\{\tiny (drop)}};
  \node[trnblk,minimum width=1.4cm, right=0.5cm of l3] (l4) {Layer 4\\{\tiny (stored)}};
  \node[grayblk,minimum width=1.4cm, right=0.5cm of l4] (l5) {Layer 5\\{\tiny (drop)}};
  \node[mrgblk,minimum width=1.4cm, right=0.5cm of l5] (out) {Output};

  \foreach \a/\b in {in/l1, l1/l2, l2/l3, l3/l4, l4/l5, l5/out}
    \draw[arr] (\a)--(\b);

  \draw[red,thick,decorate,decoration={brace,amplitude=5pt,mirror,raise=4pt}]
    (l2.south west) -- (l3.south east) node[midway,below=9pt,font=\tiny,red]{recomputed in backward};
  \draw[green!60!black,thick,decorate,decoration={brace,amplitude=5pt,raise=4pt}]
    (l1.north west) -- (l1.north east) node[midway,above=9pt,font=\tiny,green!60!black]{checkpoint};
\end{tikzpicture}
\caption{Gradient Checkpointing: only checkpoint activations stored; intermediate activations recomputed in backward.}
\end{figure}

\subsection{State Transition Table}
\begin{table}[h]\centering
\begin{tabular}{lll}\toprule
\textbf{Property} & \textbf{Before} & \textbf{After}\\\midrule
Activation memory & $O(L \cdot d)$ & $O(\sqrt{L} \cdot d)$\\
Recomputation & None & Extra forward for non-checkpoints\\
Max batch size & Limited by activations & Larger (mem freed)\\
Compute overhead & Baseline & ${\sim}30$\% more FLOPs\\
Large model feasibility & Memory OOM & Enabled\\\bottomrule
\end{tabular}
\caption{State properties before and after Gradient Checkpointing.}
\end{table}

\newpage
%% ------------------------------------------------------------
\section{Flash Attention}

\subsection{Mathematical Foundation}
Flash Attention computes attention in IO-aware tiles that fit in SRAM, avoiding materialising the full $N \times N$ attention matrix in HBM. The standard attention complexity is $O(N^2 d)$ in FLOPs and $O(N^2)$ in memory; Flash Attention maintains $O(N^2 d)$ FLOPs but reduces memory to $O(N)$:
\begin{equation}
  \text{Attention}(Q,K,V) = \text{softmax}\!\left(\frac{QK^\top}{\sqrt{d}}\right)V
\end{equation}
Computed in blocks of size $B_r \times B_c$ where $B_r = \lceil M/4d \rceil$, $M$ = SRAM size, using online softmax to accumulate numerically stable output without full matrix materialisation.

\subsection{Architecture Flow Diagram}

\begin{figure}[h]
\centering
\begin{tikzpicture}[node distance=0.7cm and 1.0cm]
  \node[datblk] (qkv) {$Q, K, V$\\(HBM)};
  \node[grayblk, right=of qkv] (tile) {Load Tile\\$B_r \times B_c$\\to SRAM};
  \node[trnblk, right=of tile] (local) {Local Softmax\\+ $V$ accumulate};
  \node[grayblk, right=of local] (write) {Write partial\\output to HBM};
  \node[mrgblk, below=0.9cm of $(tile)!0.5!(write)$] (final) {Final $O$\\(no full $N^2$ stored)};

  \draw[arr] (qkv)--(tile);
  \draw[arr] (tile)--(local);
  \draw[arr] (local)--(write);
  \draw[arr,dashed,green!60!black] (write.east) -- ++(0.4,0) |- (tile.east) node[midway,right,font=\tiny]{next tile};
  \draw[arr] (write.south) -- ++(0,-0.3) -| (final.east);

  \node[red,font=\tiny,align=center,below=0.2cm of tile] {only tile in SRAM\\not full $N^2$};
\end{tikzpicture}
\caption{Flash Attention: tiled block-wise computation avoids materialising the full attention matrix in HBM.}
\end{figure}

\subsection{State Transition Table}
\begin{table}[h]\centering
\begin{tabular}{lll}\toprule
\textbf{Property} & \textbf{Before (Standard)} & \textbf{After (Flash Attention)}\\\midrule
Attention matrix memory & $O(N^2)$ HBM & $O(N)$ (tiles in SRAM)\\
HBM read/write passes & Multiple & Minimal (one pass)\\
Compute FLOPs & $O(N^2 d)$ & $O(N^2 d)$ (same)\\
Speed (wall-clock) & Baseline & $2$--$4\times$ faster\\
Long-sequence feasibility & OOM at $N > 4096$ & $N > 100{,}000$ feasible\\\bottomrule
\end{tabular}
\caption{State properties comparing Standard and Flash Attention.}
\end{table}

\newpage
%% ------------------------------------------------------------
\section{Fully Sharded Data Parallel (FSDP)}

\subsection{Mathematical Foundation}
FSDP shards model parameters, gradients, and optimizer states across $D$ devices, reducing per-device memory from $O(|\theta|)$ to $O(|\theta|/D)$. Before each layer's forward/backward, an All-Gather reconstructs the full parameter shard:
\begin{equation}
  \theta_{\text{full}} = \text{AllGather}(\theta^{(1)}, \ldots, \theta^{(D)}),\quad \theta_{\text{full}} \in \mathbb{R}^{|\theta|}
\end{equation}
After the operation, $\theta_{\text{full}}$ is discarded. Gradient shards are reduced-scattered after the backward:
\begin{equation}
  g^{(d)} = \text{ReduceScatter}(g_{\text{full}})[d],\quad \theta^{(d)} \leftarrow \theta^{(d)} - \eta g^{(d)}
\end{equation}

\subsection{Architecture Flow Diagram}

\begin{figure}[h]
\centering
\begin{tikzpicture}[node distance=0.6cm and 1.0cm]
  \node[grayblk,minimum width=1.4cm] (d1) {Device 1\\$\theta^{(1)}$};
  \node[grayblk,minimum width=1.4cm, right=0.4cm of d1] (d2) {Device 2\\$\theta^{(2)}$};
  \node[grayblk,minimum width=0.8cm, right=0.4cm of d2] (dots) {$\cdots$};
  \node[grayblk,minimum width=1.4cm, right=0.4cm of dots] (dD) {Device $D$\\$\theta^{(D)}$};

  \node[trnblk, below=0.9cm of d2] (ag) {All-Gather\\$\theta_{\text{full}}$};
  \node[frzblk, right=1.6cm of ag] (fwd) {Layer\\Forward};
  \node[lossblk, right=of fwd] (bwd) {Backward\\$g_{\text{full}}$};
  \node[orangeblk, below=0.9cm of bwd] (rs) {Reduce-Scatter\\$g^{(d)}$};
  \node[mrgblk, left=2.2cm of rs] (upd) {Local\\Adam Update};

  \draw[arr] (d1.south) -| (ag.north);
  \draw[arr] (d2.south) -- (ag.north);
  \draw[arr] (dD.south) -| (ag.north);
  \draw[arr] (ag)--(fwd);
  \draw[arr] (fwd)--(bwd);
  \draw[arr] (bwd)--(rs);
  \draw[arr] (rs)--(upd);
  \draw[arr,dashed,green!60!black] (upd.north) |- (d2.south) node[near start,left,font=\tiny]{re-shard};
\end{tikzpicture}
\caption{FSDP: parameters sharded across devices; All-Gather for compute, Reduce-Scatter for gradient aggregation.}
\end{figure}

\subsection{State Transition Table}
\begin{table}[h]\centering
\begin{tabular}{lll}\toprule
\textbf{Property} & \textbf{Before (DDP)} & \textbf{After (FSDP)}\\\midrule
Per-device param memory & $O(|\theta|)$ & $O(|\theta|/D)$\\
Optimizer state memory & $O(|\theta|)$ per device & $O(|\theta|/D)$ per device\\
Communication pattern & AllReduce (gradients) & All-Gather + Reduce-Scatter\\
Max model size & GPU VRAM limited & $D \times$ larger\\
Training throughput & Baseline & Near-linear scaling\\\bottomrule
\end{tabular}
\caption{State properties comparing DDP and FSDP.}
\end{table}

\newpage
%% ------------------------------------------------------------
\section{DeepSpeed ZeRO}

\subsection{Mathematical Foundation}
DeepSpeed ZeRO (Zero Redundancy Optimizer) progressively eliminates redundancy across three stages. For $D$ devices and mixed-precision with Adam, the per-device memory at each stage:
\begin{align}
  \text{ZeRO-1:}&\quad \text{optimizer states sharded: } 4|\theta|/D + 6|\theta|\\
  \text{ZeRO-2:}&\quad \text{+ gradients sharded: } 4|\theta|/D + 4|\theta|/D + 4|\theta|\\
  \text{ZeRO-3:}&\quad \text{+ parameters sharded: } (4+4+4)|\theta|/D
\end{align}
ZeRO-3 achieves $12|\theta|/D$ bytes per device vs $16|\theta|$ bytes for full mixed-precision without sharding.

\subsection{Architecture Flow Diagram}

\begin{figure}[h]
\centering
\begin{tikzpicture}[node distance=0.5cm and 0.9cm]
  \node[datblk,minimum width=2.2cm] (full) {Standard\\DDP:\\$16|\theta|$ per GPU};
  \node[grayblk,minimum width=2.2cm, right=0.7cm of full] (z1) {ZeRO-1:\\Shard Opt States\\$12|\theta|+ 4|\theta|/D$};
  \node[trnblk,minimum width=2.2cm, right=0.7cm of z1] (z2) {ZeRO-2:\\+ Shard Gradients\\$4|\theta|+8|\theta|/D$};
  \node[mrgblk,minimum width=2.2cm, right=0.7cm of z2] (z3) {ZeRO-3:\\+ Shard Params\\$12|\theta|/D$};

  \draw[arr] (full)--(z1);
  \draw[arr] (z1)--(z2);
  \draw[arr] (z2)--(z3);

  \node[red,font=\tiny,below=0.25cm of full] {full redundancy};
  \node[green!60!black,font=\tiny,below=0.25cm of z3] {min per-GPU};
\end{tikzpicture}
\caption{DeepSpeed ZeRO stages: each stage shards an additional component, progressively reducing per-GPU memory.}
\end{figure}

\subsection{State Transition Table}
\begin{table}[h]\centering
\begin{tabular}{lll}\toprule
\textbf{Property} & \textbf{Standard DDP} & \textbf{ZeRO-3}\\\midrule
Optimizer state memory & $8|\theta|$ per device & $8|\theta|/D$\\
Gradient memory & $4|\theta|$ per device & $4|\theta|/D$\\
Parameter memory & $4|\theta|$ per device & $4|\theta|/D$\\
Total per device & $16|\theta|$ & $16|\theta|/D$\\
Max trainable size & GPU VRAM bound & $D \times$ larger\\\bottomrule
\end{tabular}
\caption{Memory comparison between standard DDP and DeepSpeed ZeRO-3.}
\end{table}

